\resumeSubHeadingListStart
    \resumeSubheading
    {MLOps Engineer, MLOps}{Oct 2021 -- Jun 2022}
    {}{중구, 서울}

    \hypertarget{mlops}{
    \resumeProjectHeading
    {\textbf{MLOps 기반 인스타그램 관계망 분석 및 투자 인사이트 도출 PoC}}{MLOps Engineer}
    }
    \resumeItemListStart
        \resumeItem{\textbf{핵심 문제 및 목표:} 전문가의 정성적 판단에 의존하던 아티스트의 시장 영향력을 인스타그램 관계망 데이터로 객관화/정량화. 데이터 기반으로 성장 가능성을 증명하여, \textbf{Series A 투자 유치를 성공시키는 것}을 최종 목표로 함. }
        \resumeItem{\textbf{나의 역할:} MLOps 엔지니어로서, 데이터 수집부터 \textbf{DVC를 이용한 데이터/모델 버저닝}, 여러 모델의 입력값 가공, K-Means/PCA 분석 및 결과 시각화까지 PoC의 \textbf{재현 가능한 End-to-End 파이프라인 구축을 주도}함. }
        \resumeItem{\textbf{문제 해결 과정:}}
            \resumeItemListStart
                \resumeItem{\textbf{기술적 고민 및 선택 (Why):} 당시 머신러닝 모델 자체에 대한 깊은 지식이 부족했기에, ML 리서치 엔지니어와 협업하여 각자의 전문 분야에 집중하는 방식을 선택. 동료는 \textbf{Node2Vec, GAT} 등 다양한 그래프 임베딩 모델링을, 저는 데이터 파이프라인과 후속 분석/시각화에 집중함.}
                \resumeItem{\textbf{구체적인 구현 내용 (How):} Git과 연동되는 \textbf{DVC(Data Version Control)를 도입}하여, 모델의 입력 데이터와 산출물을 Git 커밋 해시와 연동하여 버저닝하고 실험의 \textbf{재현성을 확보}함. 이후, 추출한 임베딩 벡터를 입력받아 \textbf{K-Means 클러스터링}과 \textbf{PCA}로 작가 그룹을 분석함. 최종적으로 이 분석 결과를 \textbf{Next.js 및 Cytoscape.js}를 활용하여 상호작용 가능한 네트워크 시각화 대시보드로 구현함. }
            \resumeItemListEnd
        \resumeItem{\textbf{결과 및 임팩트:}}
            \resumeItemListStart
                \resumeItem{\textbf{비즈니스 임팩트:} 개발된 시각화 대시보드의 분석 자료를 통해 아티스트의 성장 가능성을 증명, \textbf{Series A 투자 유치를 성공시키는 데 결정적인 기여}를 함. }
                \resumeItem{\textbf{개인 및 기술적 성장:} 이 협업 경험을 통해 머신러닝 파이프라인 구축에 대한 깊은 실무 이해를 얻었으며, 데이터 엔지니어로서의 커리어로 본격적으로 전환하는 결정적인 계기가 됨.}
            \resumeItemListEnd
        \hyperlink{mlops_back}{\underline{go back}}
    \resumeItemListEnd

    \hypertarget{instagram_scraper}{
    \resumeProjectHeading
    {\textbf{인스타그램 봇 탐지 우회를 위한 데이터 스크래퍼 개발 및 고도화}}{}
    }
    \resumeItemListStart
    \resumeItem{\textbf{핵심 문제 및 목표:} 인스타그램의 강력한 \textbf{봇 탐지(Bot Detection) 시스템}으로 인해 데이터의 지속적인 수집이 불가능한 문제를 해결. 봇 탐지를 안정적으로 우회하여, 대규모 소셜 관계 데이터를 \textbf{중단 없이 수집}할 수 있는 견고한 스크래퍼 개발을 목표로 함.}
    \resumeItem{\textbf{나의 역할:} 봇 탐지 문제 해결을 위한 다양한 기술적 접근법(\textbf{Selenium}, \textbf{requests}, \textbf{Instagrapi} 등)을 리서치, 프로토타이핑하고, 최종 솔루션을 구현하는 역할을 담당함.}
    \resumeItem{\textbf{문제 해결 과정:}}
        \resumeItemListStart
            \resumeItem{\textbf{트레이드오프 분석 (서버 vs. 모바일 환경):} 서버 환경에서 Python 기반의 스크래퍼를 사용하는 것은 결국 봇 탐지에 의해 차단됨을 확인함. '서버 환경' 자체가 탐지의 원인이라는 가설 하에, 초기 비용이 높더라도 \textbf{Python 라이브러리를 Swift로 직접 마이그레이션}하여 모바일 환경 요청을 흉내 내는 근본적인 해결책을 선택함.}
            \resumeItem{\textbf{구체적인 구현 내용 (How):} Python으로 작성된 \textbf{`Instagrapi`의 핵심 로직을 Swift 언어로 직접 포팅}하여, iOS 모바일 환경에서 동작하는 데이터 스크래퍼를 개발함. 개발 과정에서 발견한 라이브러리 버그를 수정하여 \textbf{오픈소스 프로젝트에 직접 기여}함.}
        \resumeItemListEnd
    \resumeItem{\textbf{결과 및 임팩트:}}
        \resumeItemListStart
            \resumeItem{\textbf{핵심 문제 해결:} 기존 방법으로는 불가능했던 \textbf{대규모 소셜 관계 데이터를 안정적으로 수집}하는 데 성공하여, 'MLOps 분석 PoC' 프로젝트의 기반을 마련함.}
            \resumeItem{\textbf{문제 해결 역량 입증:} 프레임워크를 다른 언어로 포팅하는 창의적인 접근을 통해, 고질적인 외부 시스템의 제약사항을 해결함.}
        \resumeItemListEnd
        \hyperlink{instagram_scraper_back}{\underline{go back}}
    \resumeItemListEnd

    \hypertarget{postgresql}{
    \resumeProjectHeading
    {\textbf{PostgreSQL 스키마 분리를 이용한 초기 데이터 파이프라인 구축}}{}
    }
    \resumeItemListStart
        \resumeItem{\textbf{핵심 문제 및 목표:} 수집한 원본(Raw) 데이터와 정제된(Purified) 데이터를 효율적으로 분리 및 관리할 필요가 있었음. 별도의 웨어하우스 구축이 어려운 상황에서, \textbf{단일 DB 내에서 데이터 무결성을 보존하고 안정적으로 데이터를 제공}하는 실용적인 파이프라인 구축을 목표로 함.}
        \resumeItem{\textbf{나의 역할:} 데이터 정제 로직을 개발하고, \textbf{PostgreSQL 내에서 스키마를 분리하여 데이터 계층을 나누는 아키텍처를 직접 설계하고 구현}함.}
        \resumeItem{\textbf{문제 해결 과정:}}
            \resumeItemListStart
                \resumeItem{\textbf{트레이드오프 분석 (DB 분리 vs. 스키마 분리):} 물리적인 DB 분리는 비용과 관리 포인트가 2배로 늘어나는 단점이 있었음. 대신, 추가 비용 없이 논리적 분리가 가능한 \textbf{PostgreSQL의 스키마 분리 방식}을 선택하여 실용성과 비용 효율성을 확보함.}
                \resumeItem{
                \textbf{구체적인 구현 내용 (How):} 
                \textbf{PostgreSQL} 내에 raw\_data와 purified\_data스키마를 생성. 
                \textbf{Python의 psycopg2 라이브러리}를 사용하여 데이터 정제 로직을 수행하고, 각 스키마의 테이블에 데이터를 저장하는 파이프라인을 구축함.
                }
            \resumeItemListEnd
        \resumeItem{\textbf{결과 및 임팩트:}}
            \resumeItemListStart
                \resumeItem{\textbf{데이터 무결성 확보:} 원본 데이터는 변경 없이 보존하고, 분석가는 정제된 데이터만 안전하게 사용할 수 있는 환경을 구축하여 데이터 관리의 복잡성을 낮춤.}
                \resumeItem{\textbf{성장의 발판 마련:} 이 경험을 통해 데이터 웨어하우스의 필요성을 깨닫고, 이후 더 고도화된 데이터 아키텍처를 설계하는 역량의 기반을 다짐.}
            \resumeItemListEnd
        \hyperlink{postgresql_back}{\underline{go back}}
    \resumeItemListEnd
    
\resumeSubHeadingListEnd